\chapter{General Introduction}
\label{ch:pfe-introduction}

\section{Context}

Modern software is rarely developed as an isolated body of source code.
Applications combine first-party logic with operating-system services,
open-source libraries, commercial components, build plug-ins, package
registries, and automated delivery infrastructure. This network of producers,
tools, services, and artifacts forms a software supply chain. A weakness or
compromise at any point in that chain can affect the delivered application,
which is why supply-chain risk management must cover the complete software
life cycle \cite{nist2022csc}.

Third-party components are especially important in mobile development. They
provide networking, cryptography, database access, user-interface widgets,
media handling, dependency injection, analytics, and many other capabilities.
Reuse accelerates development, but it also transfers maintenance,
licensing, and security risk into every application that embeds the component.
When a vulnerability is disclosed, an organization must determine which
applications contain the affected package and version before it can prioritize
remediation.

A Software Bill of Materials (\SBOM) supports this need by representing
software components, identifiers, versions, hashes, and relationships in a
structured form. The NTIA describes an \SBOM as a nested inventory for
software and positions it as a foundation for software transparency
\cite{ntiaSbom}. Standards such as \SPDX and \CycloneDX make this inventory
machine-readable and interoperable with governance and vulnerability tools
\cite{spdxSpecification,cyclonedxSpecification}.

The Android ecosystem creates a difficult setting for this transparency.
Applications are commonly written in Kotlin or Java, assembled with Gradle,
and distributed as compiled Android packages. External JAR and AAR artifacts
are resolved, transformed, merged with application code, converted to DEX,
combined with resources and native libraries, and signed. The resulting
\APK is the artifact installed or audited; it does not contain the original
dependency declaration in a form equivalent to the resolved build graph.

Figure~\ref{fig:pfe-intro-supply-chain} summarizes the motivating situation.
The same dependency that enables application functionality can also become a
source of inherited vulnerability or provenance risk. The security analyst
may receive only the delivered binary, after the original supply-chain
metadata has been transformed or removed.

\begin{figure}[H]
\centering
\begin{tikzpicture}[
  node distance=8mm and 9mm,
  stage/.style={draw, rounded corners, align=center, text width=0.18\textwidth,
    minimum height=12mm, fill=blue!5},
  dependency/.style={draw, rounded corners, align=center,
    text width=0.20\textwidth, minimum height=12mm, fill=warningorange!13},
  analyst/.style={draw, rounded corners, align=center,
    text width=0.19\textwidth, minimum height=12mm, fill=evidencegreen!11},
  arrow/.style={-{Latex[length=2.5mm]}, thick, shorten >=1.5pt,
    shorten <=1.5pt},
  risk/.style={-{Latex[length=2.5mm]}, dashed, thick, red!65!black,
    shorten >=1.5pt, shorten <=1.5pt}
]
\node[stage] (source) {Application source\\and build configuration};
\node[dependency, right=of source] (deps) {Third-party and\\transitive dependencies};
\node[stage, right=of deps] (build) {Build, D8/\Reight,\\merge, package, sign};
\node[stage, right=of build] (apk) {Delivered\\compiled APK};
\node[analyst, below=15mm of apk] (analyst) {Auditor or security team\\needs components and versions};
\node[dependency, below=15mm of deps] (advisory) {Vulnerability and\\provenance information};

\draw[arrow] (source) -- (deps);
\draw[arrow] (deps) -- (build);
\draw[arrow] (build) -- (apk);
\draw[arrow] (apk.south) -- (analyst.north);
\draw[risk] (advisory) -- (deps);
\draw[arrow] (advisory.east) -- (analyst.west);
\end{tikzpicture}
\caption{Android software-supply-chain context: dependencies enter before the
build, while the auditor may observe only the compiled APK.}
\label{fig:pfe-intro-supply-chain}
\end{figure}

\section{Problem Statement}

\subsection{The Source-to-Binary Information Gap}

Conventional \SBOM generation is most reliable when it runs inside a trusted
build environment. Package-manager metadata, lock files, dependency-resolution
results, source modules, and compiler information can be collected before
they disappear into the final binary. This information makes it possible to
name direct and transitive dependencies and to preserve their declared
coordinates.

That assumption does not hold in every audit scenario. A security analyst,
application store, incident-response team, customer, or third-party assessor
may possess only an \APK. The source repository may be unavailable,
incomplete, private, outdated, or inconsistent with the distributed binary.
In this case, the \SBOM must be recovered from evidence that survived inside
the compiled artifact.

The source and binary views are not equivalent. A source dependency graph may
contain build-time, metadata-only, optional, or unused dependencies that
contribute no executable code. Conversely, the \APK may include shaded,
vendored, generated, repackaged, or native code that does not correspond to a
simple direct dependency declaration. The objective is therefore not to
recreate an imaginary perfect Gradle file. It is to identify which claims can
be supported by the delivered binary and to state where the reconstruction is
incomplete.

\subsection{Loss and Transformation of Component Evidence}

During Android compilation, classes from the application and its dependencies
are converted to DEX and commonly merged into shared files. Original JAR and
AAR boundaries disappear. Large applications may use several DEX files, while
variant-aware dependency resolution can select Android-, JVM-, Kotlin-, or
other target-specific artifacts whose physical names differ from the public
dependency coordinate.

\Reight further combines shrinking, optimization, obfuscation, and bytecode
rewriting \cite{androidShrinkCode}. It may:

\begin{itemize}
  \item rename packages, classes, methods, and fields;
  \item remove unused component code;
  \item inline methods into application or library callers;
  \item merge classes or change package organization;
  \item simplify instructions and constants; and
  \item eliminate the few methods or literals that distinguish adjacent
  library versions.
\end{itemize}

These transformations affect evidence families differently. Names may vanish
while opcode or dependency patterns remain. Structural evidence may survive
renaming but change under inlining or class merging. A component identity may
be recognizable even when its exact version is not. Android
third-party-library research consequently distinguishes library detection
from the more difficult version-identification problem
\cite{backes2016libscout,zhan2021atvhunter,zhou2025sad}.

Recent common-benchmark evidence also shows why results from selected
libraries or convenient datasets cannot be treated as universal deployment
accuracy. Detector behavior changes with partial inclusion, optimization,
reference coverage, version similarity, and dataset construction
\cite{zhan2021systematic,gu2025tplbenchmark}.

\subsection{The Vulnerability-Authority Problem}

Vulnerability databases generally associate an advisory with a package
identity and an affected version range. The inference chain is:

\[
\text{APK evidence}
\rightarrow
\text{component coordinate}
\rightarrow
\text{exact version}
\rightarrow
\text{affected range}.
\]

A wrong coordinate can retrieve unrelated advisories. A wrong version can
create a false alarm or hide an affected release. Even a correct
package/version mapping does not prove runtime exploitability, because
configuration, reachability, platform behavior, and mitigations may matter.

The central problem addressed by this thesis can therefore be stated as
follows:

\begin{quote}
How can a system recover an auditable and standards-based software-component
inventory from a compiled Android APK, estimate component identities and exact
versions from incomplete or transformed evidence, and support vulnerability
analysis without converting ambiguity into false precision?
\end{quote}

Figure~\ref{fig:pfe-intro-source-binary} contrasts the two generation
conditions. The source-based path can use declared and resolved build
metadata. The APK-only path must reconstruct bounded claims from static
observations and known references.

\begin{figure}[H]
\centering
\begin{tikzpicture}[
  node distance=7mm and 11mm,
  source/.style={draw, rounded corners, align=center, text width=0.23\textwidth,
    minimum height=11mm, fill=evidencegreen!10},
  binary/.style={draw, rounded corners, align=center, text width=0.23\textwidth,
    minimum height=11mm, fill=warningorange!12},
  output/.style={draw, rounded corners, align=center, text width=0.24\textwidth,
    minimum height=11mm, fill=schoolblue!13},
  arrow/.style={-{Latex[length=2.5mm]}, thick},
  uncertain/.style={-{Latex[length=2.5mm]}, dashed, thick}
]
\node[source] (manifest) {Source and build view\\manifests, lock files, resolved graph};
\node[output, right=of manifest] (sourcebom) {Build-generated SBOM\\declared and resolved metadata};
\node[binary, below=14mm of manifest] (apk) {Binary-only view\\APK, DEX, manifest, resources, native files};
\node[binary, right=of apk] (recovery) {Evidence recovery\\reference comparison, confidence, ambiguity};
\node[output, right=of recovery] (recovered) {Recovered SBOM\\observed and explicitly bounded claims};

\draw[arrow] (manifest) -- (sourcebom);
\draw[arrow] (apk) -- (recovery);
\draw[arrow] (recovery) -- (recovered);
\draw[uncertain] (sourcebom) -- node[right, font=\scriptsize, align=left]
  {related views,\\not identical} (recovered);
\end{tikzpicture}
\caption{Source-based SBOM generation and binary-only SBOM recovery answer
related but non-identical composition questions.}
\label{fig:pfe-intro-source-binary}
\end{figure}

\section{Motivation}

\subsection{Security Auditing without Source Access}

Organizations often need to inspect software they did not build. Acquired
applications, legacy products, third-party mobile clients, incident artifacts,
and app-store packages may lack a reproducible source environment. Static
\APK analysis offers a way to recover useful evidence without executing the
application or trusting vendor-provided metadata.

A practical APK-to-\SBOM tool can help an analyst:

\begin{itemize}
  \item inventory recognizable third-party JVM/Maven components;
  \item identify strong exact-version evidence and adjacent-version
  ambiguity;
  \item obtain a standard document consumable by other tools;
  \item prioritize vulnerability investigation;
  \item compare binary observations with source dependency evidence when
  source is available; and
  \item recognize native-heavy, hidden-code, or missing-reference boundaries
  that require another analysis method.
\end{itemize}

\subsection{Transparency rather than Forced Completeness}

The usefulness of an \SBOM depends on the reliability of its identifiers.
Emitting a large inventory with unsupported exact versions may be more
dangerous than emitting a smaller inventory with clear evidence and
limitations. The project is therefore motivated by \emph{evidence quality},
not by maximizing the number of rows at any cost.

This principle has three consequences. First, library identity and exact
version are treated as separate decisions. Second, ambiguous alternatives and
reference-coverage gaps are preserved instead of hidden. Third,
vulnerability findings inherit the confidence and validation status of the
underlying component.

\subsection{An End-to-End Engineering Gap}

Existing Android library-identification research contributes important
techniques based on package signatures, class dependencies, opcodes, call
structure, optimization simulation, and learned representations. However,
detector output alone does not complete the operational workflow. A usable
system also needs:

\begin{itemize}
  \item a reproducible versioned reference database;
  \item an explicit identity/version decision policy;
  \item standard \SBOM serialization;
  \item vulnerability lookup with uncertainty propagation;
  \item source-backed validation where possible;
  \item report formats for both machines and analysts; and
  \item regression gates that preserve negative experiments and prevent
  unsafe global promotions.
\end{itemize}

The project addresses this integration gap with a working MVP whose claims are
limited by its observed evidence.

\section{Research Questions}

The thesis evaluates the following four research questions. They are phrased
to match the evidence layers and answers reported in
Chapter~\ref{ch:pfe-evaluation}.

\begin{enumerate}[label=\textbf{RQ\arabic*:}]
  \item Can the pipeline recover exact Maven versions when ground truth and
  adequate reference evidence are available?
  \item How does emitted-version quality change between clean and difficult
  real APK families?
  \item How much of the full source dependency denominator is represented by
  the reference database and emitted by the matcher?
  \item Can SBOM, vulnerability, and AI assistance preserve uncertainty and
  authority boundaries?
\end{enumerate}

These questions separate capability from coverage. RQ1 tests the mechanics
under favorable evidence. RQ2 examines realistic variation across application
families. RQ3 distinguishes reference coverage from matcher recall. RQ4
evaluates whether downstream outputs remain honest when the analysis is
uncertain.

\section{Objectives}

\subsection{General Objective}

The general objective is to design, implement, and evaluate a static
evidence-based system that recovers a bounded \CycloneDX \SBOM from a compiled
Android \APK and supports confidence-aware vulnerability analysis.

\subsection{Specific Objectives}

The specific objectives are:

\begin{enumerate}
  \item inspect an \APK and recover manifest, DEX, resource, native-file, and
  forensic metadata without installing or executing the application;
  \item extract comparable packages, classes, methods, instructions, literals,
  API references, and dependency structures from all successfully parsed DEX
  files;
  \item acquire selected Maven JAR/AAR releases, convert their bytecode to DEX,
  and construct a versioned SQLite reference-feature database;
  \item combine several static evidence families to rank component and version
  candidates;
  \item distinguish component identity from exact-version confidence and
  expose ambiguity or abstention when evidence is insufficient;
  \item emit accepted components in \CycloneDX JSON using Package URL
  identifiers and evidence properties;
  \item query ecosystem-aware vulnerability information while preserving
  component and version uncertainty;
  \item validate emitted Maven rows against source dependency trees where
  auditable source is available; and
  \item evaluate controlled, clean, partial, source-known, report-only,
  failed, and limitation evidence without merging them into one misleading
  score.
\end{enumerate}

\section{Methodology Overview}

The project followed an iterative research-and-engineering methodology. It did
not begin by selecting a single detector and treating its score as a complete
\SBOM. The work progressed through the following stages:

\begin{enumerate}
  \item \textbf{literature and standard review}: analyze Android
  third-party-library detection, obfuscation resilience, realistic
  benchmarking, \SBOM standards, and vulnerability identifiers;
  \item \textbf{scope and architecture}: define a static JVM/Maven MVP,
  choose \CycloneDX and OSV, and establish native, dynamic-code, and
  heavy-obfuscation boundaries;
  \item \textbf{reference construction}: download known Maven releases,
  convert JAR/AAR bytecode to DEX, extract comparable features, and store
  coordinate/version profiles;
  \item \textbf{APK analysis and matching}: inspect the package, aggregate
  multidex evidence, generate candidates, rank references, and separate
  identity from version confidence;
  \item \textbf{reporting and validation}: generate \SBOM, JSON, Markdown, and
  HTML outputs, query vulnerabilities, and compare emitted rows with source
  dependency trees;
  \item \textbf{controlled and real evaluation}: test clear, multidex,
  \Reight, clean real, partial real, native-heavy, and hidden-code cases in
  separate evidence lanes; and
  \item \textbf{guarded improvement}: evaluate reference growth, advanced
  feature weights, rerankers, and AI assistance with target/control gates;
  promote only changes that improve the target without regressing controls.
\end{enumerate}

Figure~\ref{fig:pfe-intro-methodology} shows the evidence-gated loop used
during implementation and optimization.

\begin{figure}[H]
\centering
\resizebox{\textwidth}{!}{%
\begin{tikzpicture}[
  node distance=7mm and 8mm,
  stage/.style={draw, rounded corners, align=center, text width=0.18\textwidth,
    minimum height=11mm, fill=blue!5},
  gate/.style={draw, diamond, aspect=2.1, align=center, inner sep=1.5mm,
    fill=warningorange!14},
  outcome/.style={draw, rounded corners, align=center, text width=0.18\textwidth,
    minimum height=11mm, fill=evidencegreen!10},
  reject/.style={draw, rounded corners, align=center, text width=0.19\textwidth,
    minimum height=11mm, fill=red!7},
  arrow/.style={-{Latex[length=2.5mm]}, thick}
]
\node[stage] (idea) {Paper-backed idea\\or observed failure};
\node[stage, right=of idea] (experiment) {Isolated implementation\\and experiment};
\node[gate, right=of experiment] (gate) {Target improves\\and controls pass?};
\node[outcome, above right=8mm and 11mm of gate] (promote)
  {Promote to calibrated\\default behavior};
\node[reject, below right=8mm and 11mm of gate] (retain)
  {Keep diagnostic,\\experiment-only, or reject};
\node[stage, right=of promote] (evidence) {Record evidence,\\decision, and limitation};

\draw[arrow] (idea) -- (experiment);
\draw[arrow] (experiment) -- (gate);
\draw[arrow] (gate) -- node[above left, font=\scriptsize] {yes} (promote);
\draw[arrow] (gate) -- node[below left, font=\scriptsize] {no} (retain);
\draw[arrow] (promote) -- (evidence);
\draw[arrow] (retain) -| (evidence);
\end{tikzpicture}
}
\caption{Fail-fast, target/control methodology used to promote or reject
matcher, reference-database, R8, and AI experiments.}
\label{fig:pfe-intro-methodology}
\end{figure}

This methodology treats failed experiments as evidence. An improvement that
fixes one target but causes false versions or components in controls is not a
successful system improvement.

\section{Thesis Organization}

The remainder of the thesis is organized as follows:

\begin{itemize}
  \item \textbf{Chapter~\ref{ch:pfe-background} -- Technical Background}
  introduces software supply-chain security, \SBOM standards, vulnerability
  terminology, Android packaging and compilation, DEX evidence, reverse
  engineering, component identification, and evidence-based recovery.

  \item \textbf{Chapter~\ref{ch:pfe-ai} -- Artificial Intelligence for
  Component Identification} presents only the ML/DL methods implemented in
  the project: the structured Random Forest, full and MVP-safe Transformers,
  the raw-opcode negative result, evaluation metrics, shortcut risks, and the
  report-only authority boundary.

  \item \textbf{Chapter~\ref{ch:pfe-sota} -- State of the Art} reviews Android
  third-party-library and version-identification approaches by technique,
  compares their capabilities and limitations, and derives the research gap
  addressed by the MVP.

  \item \textbf{Chapter~\ref{ch:pfe-design} -- Proposed Solution and System
  Design} defines requirements, architecture, actors, use cases, activities,
  sequences, evidence families, version decisions, components, deployment,
  and the reference-database model.

  \item \textbf{Chapter~\ref{ch:pfe-evaluation} -- Implementation, Evaluation,
  and Discussion} describes the implemented modules and commands, defines the
  evidence layers and metrics, reports controlled and real-APK results,
  preserves failed experiments, answers the research questions, and discusses
  validity threats.

  \item \textbf{Chapter~\ref{ch:pfe-conclusion} -- General Conclusion and
  Perspectives} synthesizes the contributions and limitations and presents
  the most defensible directions for future work.
\end{itemize}
